% Source: https://artofproblemsolving.com/community/c3h2378402_radians
% Stable identity: aops-topic-2378402-q01
An arc with central angle $\theta$ degrees has length $\dfrac{\theta}{360}\cdot2\pi r$ (it is that fraction of the circumference). With $\theta=\dfrac{360}{2\pi}$:
\[\text{arc }AC=\dfrac{360}{2\pi}\cdot\dfrac{1}{360}\cdot2\pi r=\boxed{r}.\]

\begin{asy}
usepackage("amsmath");
size(4cm);
defaultpen(fontsize(10pt));
pair O=(0,0); real r=2;
draw(circle(O,r),gray);
pair A=r*dir(0), C=r*dir(degrees(1.0));
draw(O--A,blue); draw(O--C,blue);
draw(arc(O,r,0,degrees(1.0)),red+linewidth(1.4));
draw(arc(O,0.5,0,degrees(1.0)));
dot(O); dot(A); dot(C);
label("$O$",O,W); label("$A$",A,E); label("$C$",C,NE);
label("$r$",(O+A)/2,S,blue); label("$r$",(O+C)/2,NW,blue);
label("$1$ rad",0.5*dir(degrees(1.0)/2),E,fontsize(8pt));
label("arc $=r$",r*dir(degrees(1.0)/2),E,red);
\end{asy}

So a central angle of one radian cuts off an arc exactly as long as the radius; this is the usual definition of the radian, and it is why the full angle is $2\pi$ radians.
